\documentclass[10pt]{article}
\usepackage[margin=20mm]{geometry}
\usepackage{amsmath,amssymb,mathtools,fontspec,unicode-math,xurl,hyperref}
\IfFontExistsTF{DejaVu Serif}{\setmainfont{DejaVu Serif}}{\setmainfont{Latin Modern Roman}}
\setmathfont{latinmodern-math.otf}
\setlength{\parindent}{0pt}\setlength{\parskip}{6pt}\sloppy
\title{Initial scientific review of
Voronoi conjecture for five-dimensional parallelohedra}
\author{}
\date{26 October 2026}
\begin{document}
\maketitle
Abstract: No substantive issue was established in the inspected scope. The review followed the reduction to free directions or canonical scaling, checked the local parity arguments, and examined the dependence on prior classification results. A short argument below makes explicit why a translated pair cannot leave the specified parity hyperplane in the no-free-direction branch. Exact finite checks corroborate this argument and several case tables. They do not reproduce the external classification of 110244 types or certify every cited geometric result. The reviewed object is arXiv:1906.05193v6, not an independently acquired journal version. Scientific acceptance remains pending.\par
\section*{Identity and actual coverage}
The assigned paper is Alexey Garber, Voronoi conjecture for five-dimensional parallelohedra, M25-0065. The acquired original is arXiv:1906.05193v6, dated 31 January 2025, 41 pages [paper]. Page numbers in this report are the manuscript’s printed pages, equal to one-based PDF pages. The associated journal record is Inventiones mathematicae 240 (2025), 587–631, DOI 10.1007/s00222-025-01325-0 [journal]. That record does not establish byte or mathematical equivalence with the reviewed manuscript.\par
All 1978 lines returned by the native PDF text extraction were read, including references and Appendix A. Twenty-two original PDF pages were inspected as images: \[1,4,7,11,12,14,15,16,17,19,20,23,24,26,28,29,30,31,32,33,34,35.\] This includes all five figure pages, 4, 7, 26, 30 and 33. Static inspection of the supplied v6 TeX covered selected passages concerning translated pairs, Lemma 7.2, coherence and external dependencies; the whole TeX file was not independently read or compiled. All ten manifest-listed files were hash checked.\par
The reading was systematic and claim focused. It was not a formal verification or a complete independent reconstruction of every geometric subcase and cited theorem. The absence of executable author code in this source package is not a defect.\par
\section*{Definitions, hypotheses and dependence}
The central claim, Theorem 4.1 on p. 12, concerns every convex five-dimensional polytope that tiles \(\mathbb R^5\) by translations: it is an affine image of a lattice Dirichlet–Voronoi polytope. The argument uses the associated face-to-face lattice tiling. A dual cell \(\mathcal D(F)\) is the set of centers of tiles incident to a face \(F\); its dimension label initially records the codimension of \(F\). One must not assume the general dimension conjecture for dual cells, which the manuscript explicitly leaves open beyond the known cases.\par
The parity group \(\Lambda/2\Lambda\) and the midpoint-class group \((\tfrac12\Lambda)/\Lambda\) are distinct quotients, each isomorphic to \(\mathbb F_2^5\). In the computations below, doubling identifies a midpoint class with the sum of its endpoint parity classes. Equality of such classes alone does not imply equality of real segment vectors. Lemmas 3.3–3.5, pp. 11–12, provide the geometric translation and facet-vector restrictions needed to make that passage. Lemma 3.6 excludes the parity sum of a facet triangle; Lemma 3.7 turns a covered four-dimensional midpoint hyperplane into a free direction. Its proof uses the nontrivial intersection of a two-dimensional triangle-midpoint subspace with that hyperplane and the fact that a midpoint representing a facet cannot instead represent a non-facet contact face.\par
The checked dependence chain has the following structure. Lemma 5.3 obtains the Voronoi conclusion from a free direction, using Theorem 5.2, the four-dimensional theorem, the cited combinatorially Voronoi result, and Lemma 2.15. Definition 5.1 requires equivalence of tilings together with a compatible lattice isomorphism; equality of the face lattices of the individual polytopes alone is insufficient. Theorem 5.2 constructs this stronger equivalence by separating the tiling into lattice layers and taking the segment long enough that only neighboring layers meet. The proof was checked at the level of the stated layer properties and the induced face-incidence correspondence; the external projection and saturation results were not reproved.\par
A cubical dual 3-cell produces a free direction by Lemma 6.1. In the remaining prismatic case, Lemmas 7.1–7.2 reduce the primal face to a triangle and its three edge-dual cells to prisms over tetrahedra or pyramids over triangular prisms. Their number of prismatic cells is \(3,2,1,\) or \(0\), giving the exhaustive four branches of Sections 8–11. Section 8 uses a disconnected red Venkov graph; Sections 9–10 reconstruct an impossible centrally symmetric tetrahedral contact face unless a free direction already exists. Section 11 extends gain across parallelograms and proves coherence or a free direction. The local cycle relations and their use in constructing canonical scaling were followed; the full external local Ordine argument remains a cited input.\par
\section*{A checked expansion of the translation argument}
The assertions on p. 17, Case 7.1.2 and Subcase 7.1.2.00, and the introductory paragraph of Case 7.1.3 on p. 18 restrict translated endpoints to the parity plane \(x_4=0\). A lattice translation does not preserve that plane merely because the original segment is parallel to it. The following verification supplies the relevant alternative: an escaped translated pair itself forces a free direction, so it is excluded in the branch where no free direction exists. This is a proof clarification, not a counterexample to the lemma.\par
Let \(e_1,\ldots,e_5\) be a basis of \(\mathbb F_2^5\), and put \[F=\langle e_1,e_2,e_5\rangle,\qquad
 P_0=\{0,e_1,e_2,e_5,e_1+e_5,e_2+e_5\},\qquad A=e_3.\] These are the six prism classes and one additional class in the normalization on p. 16. Write \[M(S)=\{s+t:s,t\in S\}.\] Then \(M(P_0)=F\). Suppose a translated pair, lying in the same edge-dual cell, has midpoint class encoded by \(A+v\), where \(v\in P_0\), but has endpoints outside \(x_4=0\). Their classes have the form \[R=e_4+u,\qquad S=e_3+e_4+u+v,\qquad u\in F.\] Choose \(i\in\{1,2\}\) with \(v_i=0\); this is possible for every \(v\in P_0\). Define the nonzero linear functional \[L(w)=w_i+u_iw_4,\qquad K=\ker(L|_F).\] The four-element subspace \(K\) is contained in \(P_0\): the only elements of \(F\) omitted from \(P_0\) have both first two coordinates equal to one. Moreover, \[L(A)=L(R)=L(S)=0.\] The four distinct cosets \[K,\quad A+K,\quad R+K,\quad S+K\] therefore fill \(\ker L\), a four-dimensional subspace. Each lies in \(M(P_0\cup\{A,R,S\})\): the first lies in \(M(P_0)\) and the others arise by pairing \(A,R,S\) with points of \(K\). Thus the midpoint classes of this edge-dual cell contain \(\ker L\). Lemma 3.7 gives a free direction. Consequently the endpoint restriction is justified when the proof continues under the complementary no-free-direction assumption.\par
This argument also explains why merely checking the printed parity table would be inadequate. The reviewer-written computation exhausts all \(6\cdot8=48\) choices of \((v,u)\) and finds a hyperplane witness in every case. The analytic argument establishes the reason for those witnesses; the finite run corroborates the encoding.\par
\section*{Finite checks and their mathematical limits}
The exact \(\mathbb F_2^5\) checks in `parity\_checks.py` use integer XOR, with \(e_i\) represented by bit \(i-1\). No floating-point tolerance or author program is involved. All assertions passed. In addition to the 48 cases above, the checks established:\par
- the six midpoint completions in Case 7.1.1, and the sixteen midpoint classes of a cube plus an outside parity class;\par
- the partition in Lemma 9.1, p. 24: ten occupied classes, fourteen distinct triangle-forbidden classes, and exactly eight remaining classes in \(x_3=x_4=1\);\par
- the corresponding counts in Lemma 10.1, pp. 27–28: nine occupied classes, twelve distinct forbidden classes, and eleven remaining classes;\par
- all fifteen nonzero first-four-coordinate classes required for Lemma 8.3, p. 23, using actual diagonal midpoint classes of the indicated parallelograms; the argument only requires one non-facet class in each pair differing in the fifth coordinate;\par
- the fourteen midpoint classes and the two missing classes in Case 10.1.00110;\par
- the completion used in Lemma 11.3, p. 33. Of 32 parity/orientation possibilities for the second prism, 28 would make a known facet midpoint class represent a non-facet contact face. In each of the four remaining possibilities, its two non-facet classes are precisely the two missing midpoint classes.\par
The final check explicitly uses the known facet status of the six apex edges. Counting the union of arbitrary sets of parity classes without that geometric restriction would not justify the conclusion. Conversely, these finite tests neither construct arbitrary configurations as real tilings nor prove the external classification. The geometric incidence arguments and their hypotheses were read separately from the computations.\par
\section*{Outcome, external inputs and remaining work}
No supported substantive defect was found in the inspected arguments. No claim of whole-paper correctness, independent scientific acceptance, or author adoption follows. The version record consulted on 26 October 2026 still listed v6 as the latest arXiv revision. A limited title/erratum search did not identify a specific correction; it was not an exhaustive search.\par
Section 12, pp. 34–35, explicitly identifies the 110244-type classification and the verification of a combinatorial condition on those types as inputs to Lemma 5.3. The primary classification record reports multiple implementations and a mass-formula check [classification]; that is the authors’ evidence, not a reproduction performed here. The statement of the combinatorially Voronoi dependency was cross-checked against its author record [combinatorial]. The present review did not run that classification, PORTA, or any author code. The discussion on p. 34 removes a particular Ordine computation in dimension five; it does not remove the separate classification dependency.\par
Appendix A was read as an explanation of the author’s assessment of Engel’s approach. The cited historical works were not independently audited, and no adverse finding about those works is adopted here. A further acceptance review could independently inspect the external classification certificates and the local-to-global canonical-scaling inputs. Such work concerns proof assurance, not a currently demonstrated error or a required new empirical experiment.\par
Item | Disposition\par
Issue judgment | No Issue Found within the declared scope; zero confirmed Major or Minor findings.\par
Proof clarification | The escaped-pair alternative above is analytically verified and finitely corroborated.\par
Repair/adoption | No substantive correction proposed; author adoption not requested.\par
Remaining work | Independent assessment of cited geometric and computational inputs, and separate scientific acceptance.\par
Decision | Pending; human review records remain empty.\par
9 A. Garber, Voronoi conjecture for five-dimensional parallelohedra, arXiv:1906.05193v6, 31 January 2025. Assigned PDF and supplied TeX. A. Garber, same title, Invent. Math. 240 (2025), 587–631, DOI: 10.1007/s00222-025-01325-0 (https://doi.org/10.1007/s00222-025-01325-0). Publisher metadata only. M. Dutour Sikirić, A. Garber, A. Schürmann and C. Waldmann, The complete classification of five-dimensional Dirichlet–Voronoi polyhedra of translational lattices, Acta Crystallogr. A 72 (2016), 673–683, DOI: 10.1107/S2053273316011682. Primary abstract at https://scholarworks.utrgv.edu/mss\_fac/114/; consulted only for the dependency and reported validation scope. M. Dutour Sikirić, A. Garber and A. Magazinov, On the Voronoi conjecture for combinatorially Voronoi parallelohedra in dimension five, arXiv:1812.02964v1, 7 December 2018; journal DOI: 10.1137/18M1235004. Primary abstract/statement scope consulted; not independently reproduced.\par
\section*{Supporting checks}
The supplied v6 TeX SHA-256 is\par
The actual finite-check output SHA-256 is\par
The actual run used the existing native Python interpreter on `parity\_checks.py` and exited successfully. The text-read ranges, image checks, static-source ranges and literature checks are recorded in the accompanying evidence files.\par
This report is compiled from its exact saved TeX by the authorized Tectonic 0.17.0 executable, with `–untrusted`, retained logs, intermediates and makefile dependency rules. Compiler SHA-256:\par
The receipt binds the exact source, compiler, PDF and logs. The dependency inventory is a full cache snapshot, a superset of used files. Final PDF pages are individually inspected; their image hashes and the final source/PDF hashes are recorded externally in the page-check receipt and provenance.\par
An unverified v2 HTML URL for the combinatorial dependency failed; only its v1 record was relied on. An oversized image batch was not counted as inspected; all 22 source pages were subsequently displayed successfully.\par
This assistant performed the science, calculations, static inspection, compilation and visual checks; Work only scheduled. Human reviews are empty; scientific decision pending.\par
\end{document}
